\documentclass[11pt]{article}

\usepackage[a4paper,margin=1in]{geometry}
\usepackage{amsmath,amssymb,mathtools,bm}
\usepackage{algorithm}
\usepackage{algorithmic}

\newcommand{\R}{\mathbb{R}}
\newcommand{\PiSet}{\Pi}
\newcommand{\KL}{\operatorname{KL}}
\newcommand{\diag}{\operatorname{diag}}
\newcommand{\supp}{\operatorname{supp}}

\title{Weak Optimal Transport with a Barycentric Projection Cost}
\author{}
\date{}

\begin{document}

\maketitle

\section{Continuous weak optimal transport}

Let
\[
    \alpha\in\mathcal{P}_2(\R^d),
    \qquad
    \beta\in\mathcal{P}_2(\R^d)
\]
be probability measures with finite second moments. The set of couplings between
$\alpha$ and $\beta$ is
\[
    \PiSet(\alpha,\beta)
    =
    \left\{
        \pi\in\mathcal{P}(\R^d\times\R^d):
        (\operatorname{proj}_x)_\#\pi=\alpha,
        \qquad (\operatorname{proj}_y)_\#\pi=\beta
    \right\}.
\]

Since the first marginal of $\pi$ is $\alpha$, one may disintegrate $\pi$ as
\[
    \pi(dx,dy)=\alpha(dx)\,\pi_x(dy),
\]
where $\pi_x$ is the conditional distribution of $Y$ given $X=x$. Its
conditional barycenter is
\[
    b_\pi(x)
    :=
    \int_{\R^d} y\,\pi_x(dy)
    =
    \mathbb{E}_{\pi}[Y\mid X=x].
\]

The quadratic weak optimal transport cost is
\[
    \boxed{
    \mathcal{W}_{\mathrm{weak}}^2(\alpha,\beta)
    :=
    \inf_{\pi\in\PiSet(\alpha,\beta)}
    \int_{\R^d}
    \left\|x-b_\pi(x)\right\|^2
    \,\alpha(dx).
    }
\]
Equivalently, introducing the weak cost
\[
    C(x,p)
    =
    \left\|
        x-\int_{\R^d}y\,p(dy)
    \right\|^2,
\]
one has
\[
    \mathcal{W}_{\mathrm{weak}}^2(\alpha,\beta)
    =
    \inf_{\pi\in\PiSet(\alpha,\beta)}
    \int_{\R^d} C(x,\pi_x)\,\alpha(dx).
\]

This differs from classical quadratic optimal transport:
\[
    W_2^2(\alpha,\beta)
    =
    \inf_{\pi\in\PiSet(\alpha,\beta)}
    \int_{\R^d\times\R^d}
    \|x-y\|^2\,\pi(dx,dy).
\]
Indeed, conditional Jensen's inequality gives
\[
    \left\|x-\mathbb{E}[Y\mid X=x]\right\|^2
    \leq
    \mathbb{E}\left[\|x-Y\|^2\mid X=x\right],
\]
and consequently
\[
    \mathcal{W}_{\mathrm{weak}}^2(\alpha,\beta)
    \leq W_2^2(\alpha,\beta).
\]

Moreover,
\[
    \mathcal{W}_{\mathrm{weak}}^2(\alpha,\beta)=0
\]
whenever there exists a martingale coupling satisfying
\[
    \mathbb{E}[Y\mid X]=X.
\]
In particular, by Strassen's theorem, this happens when $\alpha$ is dominated
by $\beta$ in convex order.

\section{Discrete formulation}

Assume that
\[
    \alpha=\sum_{i=1}^n a_i\delta_{x_i},
    \qquad
    \beta=\sum_{j=1}^m b_j\delta_{y_j},
\]
where
\[
    a_i>0,\qquad b_j>0,\qquad
    \sum_{i=1}^n a_i=\sum_{j=1}^m b_j=1.
\]

A discrete coupling is represented by a nonnegative transport matrix
\[
    P=(P_{ij})\in\R_+^{n\times m}.
\]
The marginal constraints are
\[
    P\bm{1}_m=\bm{a},
    \qquad
    P^\top\bm{1}_n=\bm{b},
\]
so the discrete transport polytope is
\[
    U(\bm{a},\bm{b})
    :=
    \left\{
        P\in\R_+^{n\times m}:
        P\bm{1}_m=\bm{a},\
        P^\top\bm{1}_n=\bm{b}
    \right\}.
\]

The conditional probability of $y_j$ given $x_i$ is
\[
    \mathbb{P}(Y=y_j\mid X=x_i)
    =
    \frac{P_{ij}}{a_i}.
\]
Thus, the barycentric projection of row $i$ is
\[
    \overline{y}_i(P)
    :=
    \mathbb{E}[Y\mid X=x_i]
    =
    \frac{1}{a_i}\sum_{j=1}^m P_{ij}y_j.
\]

The discrete weak transport problem is therefore
\[
    \boxed{
    \min_{P\in U(\bm{a},\bm{b})}
    \mathcal{B}(P),
    \qquad
    \mathcal{B}(P)
    :=
    \sum_{i=1}^n
    a_i
    \left\|
        x_i-\frac{1}{a_i}\sum_{j=1}^m P_{ij}y_j
    \right\|^2.
    }
\]
An equivalent expression is
\[
    \mathcal{B}(P)
    =
    \sum_{i=1}^n
    \frac{1}{a_i}
    \left\|
        a_i x_i-\sum_{j=1}^m P_{ij}y_j
    \right\|^2.
\]

Notice that the objective depends on each row of $P$ only through its
barycenter. Consequently, two different transport matrices can have exactly
the same weak transport cost.

\section{Convexity}

For fixed positive weights $a_i$, the map
\[
    P
    \longmapsto
    \frac{1}{a_i}
    \left\|
        a_i x_i-\sum_j P_{ij}y_j
    \right\|^2
\]
is a convex quadratic function of the $i$-th row of $P$. Hence
$\mathcal{B}$ is convex.

More explicitly, for a perturbation $H\in\R^{n\times m}$,
\[
    D^2\mathcal{B}(P)[H,H]
    =
    2\sum_{i=1}^n
    \frac{1}{a_i}
    \left\|
        \sum_{j=1}^m H_{ij}y_j
    \right\|^2
    \geq 0.
\]
Since $U(\bm{a},\bm{b})$ is convex and defined by linear constraints, the
discrete weak transport problem is a convex optimization problem.

It is generally not strictly convex. Indeed, if
\[
    \sum_{j=1}^m H_{ij}y_j=0
    \qquad\text{for every }i,
\]
then the quadratic form above vanishes, even when $H\neq 0$. Thus, different
couplings with identical row barycenters may all be minimizers.

The same convexity principle holds in the continuous setting because
\[
    p\longmapsto
    \left\|x-\int y\,p(dy)\right\|^2
\]
is convex in the conditional measure $p$.

\section{KL-regularized weak transport}

A natural entropic regularization uses the independent coupling
$\alpha\otimes\beta$ as reference:
\[
    \boxed{
    \inf_{\pi\in\PiSet(\alpha,\beta)}
    \left\{
        \int
        \left\|x-b_\pi(x)\right\|^2\,\alpha(dx)
        +
        \varepsilon
        \KL\!\left(\pi\,\middle\|\,\alpha\otimes\beta\right)
    \right\}.
    }
\]

In the discrete setting, define
\[
    Q_{ij}:=a_i b_j
\]
and use the generalized KL divergence
\[
    \KL(P\|Q)
    =
    \sum_{i=1}^n\sum_{j=1}^m
    \left[
        P_{ij}\log\left(\frac{P_{ij}}{Q_{ij}}\right)
        -P_{ij}+Q_{ij}
    \right].
\]
The regularized problem is
\[
    \boxed{
    \min_{P\in U(\bm{a},\bm{b})}
    \left\{
        \mathcal{B}(P)
        +
        \varepsilon\KL(P\|\bm{a}\bm{b}^{\top})
    \right\}.
    }
\]
For $\varepsilon>0$, the KL term is strictly convex on positive matrices.
Thus, under the assumption $a_i,b_j>0$, the regularized problem has a unique
optimal transport matrix.

\section{Optimality conditions}

Let
\[
    \overline{y}_i(P)
    =
    \frac{1}{a_i}\sum_j P_{ij}y_j.
\]
The gradient of the barycentric cost is
\[
    \frac{\partial\mathcal{B}}{\partial P_{ij}}(P)
    =
    2\left\langle
        \overline{y}_i(P)-x_i,\,
        y_j
    \right\rangle.
\]
Define the endogenous cost matrix
\[
    C_{ij}(P)
    :=
    2\left\langle
        \overline{y}_i(P)-x_i,\,
        y_j
    \right\rangle.
\]
The first-order optimality conditions imply that there exist potentials
$f_i$ and $g_j$ such that
\[
    P_{ij}
    =
    a_i b_j
    \exp\left(
        \frac{
            f_i+g_j-C_{ij}(P)
        }{\varepsilon}
    \right).
\]
Unlike classical entropic optimal transport, the cost matrix depends on the
unknown coupling through its row barycenters. Therefore, one does not obtain
the solution from a single Sinkhorn scaling. Instead, the solution satisfies
a nonlinear Sinkhorn fixed-point equation.

\section{Saddle-point reformulation}

Use the Fenchel identity
\[
    \|x-z\|^2
    =
    \sup_{u\in\R^d}
    \left\{
        2\langle u,x-z\rangle-\|u\|^2
    \right\}.
\]
Applying it row by row yields
\begin{align*}
    \mathcal{B}(P)
    &=
    \sup_{u_1,\ldots,u_n\in\R^d}
    \left\{
        \sum_{i=1}^n
        \left(
            2a_i\langle u_i,x_i\rangle
            -a_i\|u_i\|^2
        \right)
        -
        2\sum_{i,j}P_{ij}\langle u_i,y_j\rangle
    \right\}.
\end{align*}
Consequently, the regularized weak transport problem is equivalent to
\[
    \sup_{\bm{u}}
    \left\{
        \sum_i
        \left(
            2a_i\langle u_i,x_i\rangle-a_i\|u_i\|^2
        \right)
        +
        \min_{P\in U(\bm{a},\bm{b})}
        \left[
            \sum_{i,j}C_{ij}(\bm{u})P_{ij}
            +
            \varepsilon\KL(P\|\bm{a}\bm{b}^{\top})
        \right]
    \right\},
\]
where
\[
    C_{ij}(\bm{u})
    =
    -2\langle u_i,y_j\rangle.
\]
For fixed $\bm{u}$, the inner problem is now an ordinary entropic linear
optimal transport problem and can be solved by Sinkhorn iterations.

The corresponding Gibbs kernel is
\[
    K_{ij}(\bm{u})
    =
    \exp\left(
        -\frac{C_{ij}(\bm{u})}{\varepsilon}
    \right)
    =
    \exp\left(
        \frac{2\langle u_i,y_j\rangle}{\varepsilon}
    \right).
\]
The optimal coupling for fixed $\bm{u}$ has the scaling form
\[
    P_{ij}
    =
    a_i b_j\,r_i K_{ij}(\bm{u})s_j,
\]
where $r_i,s_j>0$ enforce the marginal constraints.

The Sinkhorn scaling iterations are
\[
    r_i
    \leftarrow
    \frac{1}{
        \sum_{j=1}^m b_j K_{ij}(\bm{u})s_j
    },
    \qquad i=1,\ldots,n,
\]
followed by
\[
    s_j
    \leftarrow
    \frac{1}{
        \sum_{i=1}^n a_i r_i K_{ij}(\bm{u})
    },
    \qquad j=1,\ldots,m.
\]

Let $P(\bm{u})$ denote the result of this inner Sinkhorn problem, and let
\[
    \overline{y}_i(\bm{u})
    :=
    \frac{1}{a_i}\sum_j P_{ij}(\bm{u})y_j.
\]
By the envelope theorem, the gradient of the outer concave objective
$\Phi(\bm{u})$ is
\[
    \nabla_{u_i}\Phi(\bm{u})
    =
    2a_i
    \left(
        x_i-u_i-\overline{y}_i(\bm{u})
    \right).
\]
At the optimum,
\[
    \boxed{
    u_i^\star
    =
    x_i-\overline{y}_i(P^\star).
    }
\]
Substituting this identity into the Gibbs kernel recovers the nonlinear
optimality condition
\[
    K_{ij}^\star
    =
    \exp\left(
        \frac{
            2\langle
                x_i-\overline{y}_i(P^\star),y_j
            \rangle
        }{\varepsilon}
    \right).
\]

\section{Sinkhorn-type algorithm}

A practical algorithm alternates between a Sinkhorn solve for fixed
$\bm{u}$ and a damped ascent step in $\bm{u}$.

\begin{algorithm}[h]
\caption{Damped Sinkhorn iteration for weak entropic transport}
\begin{algorithmic}[1]
\STATE Initialize $u_i^{(0)}=0$ and $s_j^{(0)}=1$.
\FOR{$k=0,1,2,\ldots$}
    \STATE Form the kernel
    \[
        K_{ij}^{(k)}
        =
        \exp\left(
            \frac{2\langle u_i^{(k)},y_j\rangle}{\varepsilon}
        \right).
    \]
    \STATE Using warm starts, iterate
    \[
        r_i
        \leftarrow
        \frac{1}{\sum_j b_jK_{ij}^{(k)}s_j},
        \qquad
        s_j
        \leftarrow
        \frac{1}{\sum_i a_ir_iK_{ij}^{(k)}}
    \]
    until the marginal errors are sufficiently small.
    \STATE Set
    \[
        P_{ij}^{(k)}
        =
        a_i b_jr_iK_{ij}^{(k)}s_j.
    \]
    \STATE Compute the barycentric projections
    \[
        \overline{y}_i^{(k)}
        =
        \frac{1}{a_i}
        \sum_j P_{ij}^{(k)}y_j.
    \]
    \STATE Update, with a damping parameter $\rho_k>0$,
    \[
        u_i^{(k+1)}
        =
        u_i^{(k)}
        +
        \rho_k
        \left(
            x_i-u_i^{(k)}-\overline{y}_i^{(k)}
        \right).
    \]
\ENDFOR
\end{algorithmic}
\end{algorithm}

The update is a diagonally preconditioned gradient-ascent step for the
concave outer objective. The undamped fixed-point choice $\rho_k=1$ gives
\[
    u_i^{(k+1)}
    =
    x_i-\overline{y}_i^{(k)},
\]
but this full update can oscillate. A smaller step size, or a backtracking
line search on the outer objective, gives a more robust method.

A suitable stationarity residual is
\[
    \max_i
    \left\|
        u_i^{(k)}
        -
        \left(x_i-\overline{y}_i^{(k)}\right)
    \right\|.
\]
It should be monitored together with the marginal errors
\[
    \left\|P^{(k)}\bm{1}_m-\bm{a}\right\|,
    \qquad
    \left\|(P^{(k)})^\top\bm{1}_n-\bm{b}\right\|.
\]

For small $\varepsilon$, the entries of $K$ can overflow or underflow.
As in standard entropic optimal transport, one should then implement the
Sinkhorn iterations in the logarithmic domain using log-sum-exp
stabilization.

\section{Summary}

The main conclusions are:
\begin{itemize}
    \item The barycentric weak transport objective is convex in the coupling
    $P$, because it is a sum of convex quadratic functions of row
    barycenters.
    \item It is generally not strictly convex, since the objective cannot
    distinguish couplings having the same conditional barycenters.
    \item Adding $\varepsilon\KL(P\|\bm{a}\bm{b}^{\top})$ makes the discrete
    objective strictly convex and selects a unique coupling.
    \item The regularized problem is not a standard entropic OT problem with
    a fixed cost matrix: its effective cost depends on $P$.
    \item A Fenchel saddle-point representation introduces variables $u_i$
    for which the coupling subproblem becomes an ordinary entropic transport
    problem. It can therefore be solved by an inner Sinkhorn scaling,
    combined with an outer damped ascent or fixed-point iteration.
\end{itemize}

\section{Linearization and Sinkhorn iteration}

Let
\[
    b_\pi(x):=\int_{\mathbb{R}^d}y\,\pi_x(dy)
\]
be the barycentric projection of a coupling
\(\pi(dx,dy)=\alpha(dx)\pi_x(dy)\). Define the coupling-dependent cost
\[
    \boxed{
    c_\pi(x,y)
    :=
    2\left\langle x-b_\pi(x),\,x-y\right\rangle.
    }
\]
Then the quadratic weak transport functional
\[
    \mathcal B(\pi)
    =
    \int_{\mathbb{R}^d}
    \|x-b_\pi(x)\|^2\,\alpha(dx)
\]
satisfies the exact identity
\begin{align*}
    \frac12\langle c_\pi,\pi\rangle
    &=
    \frac12
    \int_{\mathbb{R}^d\times\mathbb{R}^d}
    c_\pi(x,y)\,\pi(dx,dy)\\
    &=
    \int_{\mathbb{R}^d}
    \left\langle
        x-b_\pi(x),
        x-\int y\,\pi_x(dy)
    \right\rangle
    \alpha(dx)\\
    &=
    \int_{\mathbb{R}^d}
    \|x-b_\pi(x)\|^2\,\alpha(dx).
\end{align*}
Therefore,
\[
    \boxed{
    \mathcal B(\pi)
    =
    \mathrm{cst}
    +
    \frac12\langle c_\pi,\pi\rangle,
    \qquad
    \mathrm{cst}=0.
    }
\]

The actual first-variation cost is
\[
    \ell_\pi(x,y)
    =
    2\left\langle b_\pi(x)-x,y\right\rangle.
\]
The two costs differ only by a function of \(x\):
\[
    c_\pi(x,y)
    =
    \ell_\pi(x,y)
    +
    2\left\langle x-b_\pi(x),x\right\rangle.
\]
Since every admissible coupling has first marginal \(\alpha\), adding a
function of \(x\) to an OT cost does not change the minimizing coupling.
Thus, \(c_\pi\) can be used as the linearized OT cost.

More precisely, for any
\(\gamma\in\Pi(\alpha,\beta)\),
\[
    \mathcal B(\gamma)
    =
    \mathcal B(\pi)
    +
    \langle c_\pi,\gamma-\pi\rangle
    +
    \int
    \|b_\gamma(x)-b_\pi(x)\|^2\,\alpha(dx).
\]
Hence the first-order approximation around \(\pi\) is
\[
    \mathcal B(\gamma)
    \simeq
    \mathcal B(\pi)
    +
    \langle c_\pi,\gamma-\pi\rangle.
\]

\subsection{Discrete expression}

For
\[
    \alpha=\sum_{i=1}^n a_i\delta_{x_i},
    \qquad
    \beta=\sum_{j=1}^m b_j\delta_{y_j},
\]
let \(P\in U(\bm a,\bm b)\) and define
\[
    \overline y_i(P)
    =
    \frac{1}{a_i}\sum_{j=1}^mP_{ij}y_j.
\]
The discrete coupling-dependent cost matrix is
\[
    \boxed{
    C_{ij}(P)
    =
    2\left\langle
        x_i-\overline y_i(P),
        x_i-y_j
    \right\rangle.
    }
\]
It satisfies
\begin{align*}
    \frac12\langle C(P),P\rangle
    &=
    \frac12\sum_{i,j}P_{ij}C_{ij}(P)\\
    &=
    \sum_i a_i
    \left\langle
        x_i-\overline y_i(P),
        x_i-\overline y_i(P)
    \right\rangle\\
    &=
    \sum_i a_i
    \|x_i-\overline y_i(P)\|^2.
\end{align*}
Consequently,
\[
    \boxed{
    \mathcal B(P)
    =
    \frac12\langle C(P),P\rangle.
    }
\]

The entries can equivalently be written as
\[
    C_{ij}(P)
    =
    2
    \left\langle
        \frac1{a_i}\sum_{\ell=1}^m
        P_{i\ell}(y_\ell-x_i),
        y_j-x_i
    \right\rangle.
\]

\subsection{Linearized entropic OT subproblem}

Given the current coupling \(P^{(k)}\), compute
\[
    C_{ij}^{(k)}
    :=
    C_{ij}(P^{(k)}).
\]
The linearized KL-regularized problem is the classical entropic OT problem
\[
    \boxed{
    Q^{(k)}
    =
    \operatorname*{argmin}_{Q\in U(\bm a,\bm b)}
    \left\{
        \langle C^{(k)},Q\rangle
        +
        \varepsilon
        \KL\left(Q\,\middle\|\,\bm a\bm b^\top\right)
    \right\}.
    }
\]
Introduce the Gibbs kernel
\[
    K_{ij}^{(k)}
    =
    \exp\left(
        -\frac{C_{ij}^{(k)}}{\varepsilon}
    \right).
\]
The solution has the scaling form
\[
    Q_{ij}^{(k)}
    =
    a_i b_j
    u_i^{(k)}K_{ij}^{(k)}v_j^{(k)}.
\]
The scaling vectors are computed with standard Sinkhorn iterations:
\[
    u_i
    \leftarrow
    \frac{1}{
        \sum_{j=1}^m b_jK_{ij}^{(k)}v_j
    },
    \qquad
    v_j
    \leftarrow
    \frac{1}{
        \sum_{i=1}^n a_i u_iK_{ij}^{(k)}
    }.
\]

\subsection{Linearize-and-Sinkhorn algorithm}

A simple fixed-point iteration consists of setting
\[
    P^{(k+1)}=Q^{(k)}.
\]
A more stable version uses a relaxed update
\[
    \boxed{
    P^{(k+1)}
    =
    (1-\tau_k)P^{(k)}
    +
    \tau_k Q^{(k)},
    \qquad 0<\tau_k\leq 1.
    }
\]

The complete iteration is therefore:
\[
\boxed{
\begin{array}{ll}
\text{1.} &
\displaystyle
\overline y_i^{(k)}
=
\frac1{a_i}\sum_jP_{ij}^{(k)}y_j,
\\[1.2em]
\text{2.} &
\displaystyle
C_{ij}^{(k)}
=
2\left\langle
x_i-\overline y_i^{(k)},x_i-y_j
\right\rangle,
\\[1.2em]
\text{3.} &
\displaystyle
Q^{(k)}
=
\operatorname{Sinkhorn}
\left(
\bm a,\bm b,
\exp(-C^{(k)}/\varepsilon)
\right),
\\[1.2em]
\text{4.} &
\displaystyle
P^{(k+1)}
=
(1-\tau_k)P^{(k)}
+
\tau_kQ^{(k)}.
\end{array}
}
\]

This is an entropic generalized conditional-gradient algorithm. The
quadratic barycentric functional is linearized at the current coupling,
and the resulting linear entropic OT problem is solved by Sinkhorn.

The full step \(\tau_k=1\) gives the simplest fixed-point algorithm, but it
may oscillate. Damping, for example
\[
    \tau_k=\frac{2}{k+2},
\]
or an objective-based line search generally gives a more robust iteration.


\end{document}
